\section{Tích chập, stride, padding và kích hoạt}
\subsection{Tích chập nhiều kênh}
Xét đầu vào $X$ kích thước $C_{\rm in}\times H\times W$ và đầu ra $Y$ có $C_{\rm out}$ kênh. Mỗi kênh đầu ra $o$ có trọng số $W_{o,c,u,v}$ kích thước $C_{\rm in}\times K_h\times K_w$ và bias $b_o$. Với stride $s_h,s_w$ và padding trên/trái $p_t,p_l$, kết quả trước kích hoạt là
\begin{equation}
 Y_{o,y,x}=b_o+\sum_{c=0}^{C_{\rm in}-1}\sum_{u=0}^{K_h-1}\sum_{v=0}^{K_w-1}
 W_{o,c,u,v}X_{c,s_hy+u-p_t,s_wx+v-p_l}.
\end{equation}
Đây là tích chập nhiều kênh: mỗi đầu ra tích lũy trên toàn bộ kênh đầu vào và vị trí kernel. Trường hợp một kênh có $C_{\rm in}=1$. Theo quy ước CNN, kernel không bị lật chỉ số; phép tính về hình thức là tương quan chéo, thường vẫn gọi là tích chập.

Stride là khoảng dịch cửa sổ giữa hai vị trí đầu ra liên tiếp. Padding bổ sung giá trị ngoài biên, thường bằng không trong miền thực, và có thể không đối xứng. Với padding dưới/phải $p_b,p_r$, không dùng dilation, kích thước đầu ra là
\begin{equation}
\begin{split}
 H_{\rm out}&=\left\lfloor\frac{H+p_t+p_b-K_h}{s_h}\right\rfloor+1,\\
 W_{\rm out}&=\left\lfloor\frac{W+p_l+p_r-K_w}{s_w}\right\rfloor+1.
\end{split}
\end{equation}
Stride lớn hơn một thường giảm độ phân giải. Padding cùng stride quyết định cách xử lý cửa sổ biên và vị trí lấy mẫu.

Một MAC (multiply--accumulate) nhân hai toán hạng rồi cộng vào tổng. Lớp trên cần khoảng $H_{\rm out}W_{\rm out}C_{\rm out}K_hK_wC_{\rm in}$ MAC, chưa tính bias và hậu xử lý. Đây là khối lượng số học, không bao gồm tải dữ liệu hoặc điều khiển.

\subsection{Kích hoạt và clamp}
Hàm kích hoạt tạo tính phi tuyến. ReLU là $\max(0,x)$; ReLU6 là $\min(6,\max(0,x))$. Khi lượng tử hóa, các cận thực phải ánh xạ sang miền nguyên tương ứng. Clamp giới hạn giá trị vào $[q_{\min},q_{\max}]$. Clamp theo toàn miền INT8 chỉ giới hạn khả năng biểu diễn, không tự tương đương ReLU hoặc ReLU6. Cần phân biệt cận kiểu dữ liệu với cận kích hoạt của mô hình.

\section{Tích chập depthwise, pointwise và tách theo chiều sâu}
\subsection{Depthwise: xử lý không gian theo từng kênh}
Với depth multiplier bằng một, depthwise convolution (DW) dùng một bộ lọc không gian riêng cho mỗi kênh đầu vào và giữ nguyên số kênh:
\begin{equation}
 D_{c,y,x}=b^{\rm DW}_c+\sum_{u=0}^{K_h-1}\sum_{v=0}^{K_w-1}
 W^{\rm DW}_{c,u,v}X_{c,s_hy+u-p_t,s_wx+v-p_l}.
\end{equation}
Các kênh được xử lý độc lập, không trộn thông tin. Stride và padding tác động lên cửa sổ như tích chập chuẩn; stride~2 có thể giảm kích thước không gian mà giữ số kênh trong trường hợp này.

\subsection{Pointwise: tích vô hướng theo chiều kênh}
Pointwise convolution (PW) là tích chập $1\times1$. Với stride~1 và không padding, tại mỗi vị trí, vector $C_{\rm in}$ phần tử được nhân với vector trọng số của từng kênh đầu ra:
\begin{equation}
 P_{o,y,x}=b^{\rm PW}_o+\sum_{c=0}^{C_{\rm in}-1}W^{\rm PW}_{o,c}D_{c,y,x}.
\end{equation}
Tích vô hướng này tạo $C_{\rm out}$ kênh, giữ kích thước không gian. Cùng trọng số được dùng lại ở mọi vị trí, tạo cơ hội giữ trọng số cục bộ.

\subsection{Chi phí của cặp DW--PW}
Depthwise-separable convolution thực hiện DW rồi PW: lọc không gian trước, trộn kênh sau. Đây là thành phần nền tảng của MobileNets~\cite{howard2017mobilenets}. Với kernel $K\times K$, depth multiplier bằng một và cùng miền đầu ra, số MAC mỗi vị trí là $K^2C_{\rm in}+C_{\rm in}C_{\rm out}$, so với $K^2C_{\rm in}C_{\rm out}$ của tích chập chuẩn. Tỷ số là
\begin{equation}
 \frac{\mathrm{MAC}_{\rm DW+PW}}{\mathrm{MAC}_{\rm dense}}=\frac{1}{C_{\rm out}}+\frac{1}{K^2}.
\end{equation}
So sánh bỏ qua bias, kích hoạt và truyền dữ liệu. Cách tách thay đổi cấu trúc tham số, không bảo đảm thay thế mọi kernel chuẩn mà giữ nguyên hàm. TinyMobileNet trong luận văn dùng các cặp DW--PW tùy biến, không phải MobileNetV2.

\section{GAP, lớp kết nối đầy đủ và quyết định phân loại}
Pooling thu gọn miền không gian. MaxPool chọn phần tử lớn nhất trong cửa sổ cục bộ và chỉ được nêu làm ngữ cảnh chung. Mạng của luận văn dùng Global Average Pooling (GAP), lấy trung bình toàn bản đồ từng kênh:
\begin{equation}
 g_c=\frac{1}{HW}\sum_{y=0}^{H-1}\sum_{x=0}^{W-1}X_{c,y,x}.
\end{equation}
GAP chuyển tensor $C\times H\times W$ thành vector $C$ phần tử. Trong miền nguyên, phép chia và ánh xạ sang kiểu đầu ra cần quy tắc làm tròn và tái lượng tử hóa.

Lớp kết nối đầy đủ (fully connected, FC) ánh xạ vector đặc trưng sang điểm số của $J$ lớp:
\begin{equation}
 s_j=b_j+\sum_{c=0}^{C-1}W^{\rm FC}_{j,c}g_c,
 \qquad \widehat{j}=\operatorname*{arg\,max}_{0\leq j<J}s_j.
\end{equation}
Chuỗi GAP $\rightarrow$ FC $\rightarrow$ điểm số $\rightarrow$ argmax tạo quyết định phân loại. Softmax có thể chuyển logits thực thành phân bố chuẩn hóa, nhưng đường suy luận của luận văn chọn argmax trực tiếp. Điểm số nguyên trước Softmax không phải xác suất. Các điểm số so sánh trực tiếp cần có thang đo và cách biểu diễn tương thích.
